
Evaluating the quality of text generated by language models or written by humans has long been a challenging endeavor, consistently garnering substantial attention~\citep{celikyilmazEvaluationText2021a}. Traditional methodologies predominantly rely on human annotation of texts~\citep{callison2009fast}, an approach considered overly demanding in terms of time and cost. Automatic evaluation metrics based on n-grams, such as Rouge~\citep{lin2004rouge}, BLEU~\citep{papineni2002bleu}, and METEOR~\citep{banerjee2005meteor}, have been proposed to tackle this issue~\citep{kondrak2005n}. However, these methods have been shown to exhibit a relatively weak correlation with human judgments, particularly in the context of tasks involving open-ended generation or requiring domain-specific expertise~\citep{novikova-etal-2017-need}.

Recent advancements in the field of natural language processing have led to the emergence of billion-parameter scale LLMs, such as GPT-3~\citep{brown2020language}. These LLMs have demonstrated remarkable capabilities across diverse downstream tasks, presenting new opportunities for text quality evaluation using such models. Moreover, various training paradigms have been proposed to endow LLMs with the ability to accomplish tasks in a zero-shot manner and better adhere to human-provided instructions~\citep{ouyang2022training, sanh2021multitask, wei2021finetuned}. These advancements facilitate the prompting of LLMs to evaluate generated text, effectively simulating human evaluators in the assessment process.


In view of the impressive text understanding and instruction-following capabilities of recent LLMs, a body of literature~\citep{liu2023gpteval, chiang2023can, gao2023human, shen2023large} has adopted LLM as an evaluator to assess the quality of responses to open-ended questions or traditional NLG tasks, including dialogue response generation and summarization. This methodology is dubbed LLM-as-a-judge~\citep{zheng2023judging}. Findings from these researches indicate that LLM can mimic human behavior and provide evaluations that correspond with human judgments, revealing a potentially scalable and transparent alternative to costly and laborious human evaluations.

While a \textit{single} powerful LLM can already tackle various missions, emerging studies suggest that \textit{multiple} LLMs can further improve one another through debate and cooperation~\citep{li2023camel, liang2023encouraging}. By incorporating multiple LLMs into an integrated group and designing specific interaction mechanisms, different LLMs can engage in proposing and deliberating unique responses and thought processes across several rounds. This approach leads to enhanced factuality of generated responses~\citep{du2023improving} and improvement in the completion of arduous tasks~\citep{li2023camel, qian2023communicative}. Furthermore, the multi-agent group also addresses and mitigates the Degeneration-of-Thought (DOT) problem~\citep{liang2023encouraging}.


In the human evaluation processes, relying on a single perspective can introduce bias and instability in the results~\citep{karpinska2021perils}. Recognizing this, best practices often involve multiple human annotators collaborating in the evaluation~\citep{van2019best}. Drawing inspiration from this collaborative and iterative human evaluation approach, we propose ChatEval, a system that enables each agent to employ varied communication strategies in collaborative discussion, working towards formulating final judgments. Furthermore, to enrich the evaluation dynamics, every agent within ChatEval is endowed with a unique persona. This deliberate design ensures that each agent focuses on distinct perspectives or brings specific expertise to the table. By doing so, the collective evaluation benefits from a more comprehensive lens, capturing nuances and subtleties that a single perspective might overlook. We derive this idea primarily from the insight of \textit{'There are a thousand Hamlets in a thousand people's eyes'}, meaning that every person has their unique interpretation or perspective, especially applicable to text evaluation. Indeed, these divergent perspectives shape the comprehensive and multifaceted assessment of \textit{Hamlet}. Another underlying intuition of our work stems from renowned concepts in sociology and biology, including \textit{Collective Intelligence}\citep{woolley2010evidence} and \textit{Cognitive Synergy}\citep{luppi2022synergistic}, where multiple cognitive processes or systems interact and cooperate in a way that produces a combined effect greater than the sum of their separate effects.

To summarize, the main contribution of our work is as follows:

\begin{enumerate}
\item We propose a multi-agent-based framework called \textbf{ChatEval} that aligns better with human preferences compared with single-agent-based approaches as depicted in Figure~\ref{fig:better_compare}. 
\item We propose various communication strategies and demonstrate the necessity of diverse role prompts in multi-agent debate scenarios.
\item We release our library. It's designed to be both composable and scalable, enabling researchers to implement their unique communication strategies easily. We hope this contributes to advancing research in the field of communicative agents and beyond.
\end{enumerate}


\begin{figure}[t]
\begin{center}
\includegraphics[width=\textwidth]{figs/better_compare.pdf}
\end{center}
\caption{When several referees participate in the evaluation process, they can discuss with each other and finally give a judgment that is better aligned with human annotators.}
\label{fig:better_compare}
\end{figure}